\documentclass[11pt]{article}

\usepackage[utf8]{inputenc}
\usepackage[T1]{fontenc}
\usepackage[margin=1in]{geometry}
\usepackage{xcolor}
\usepackage{listings}
\usepackage{booktabs}
\usepackage{fancyhdr}
\usepackage{titlesec}
\usepackage{enumitem}
\usepackage[colorlinks=true,linkcolor=black,urlcolor=blue!60!black]{hyperref}

% ---- colours -----------------------------------------------------------
\definecolor{ink}{rgb}{0.12,0.12,0.14}
\definecolor{panel}{rgb}{0.96,0.96,0.97}
\definecolor{rule}{rgb}{0.55,0.55,0.60}

% ---- calculator I/O listing style --------------------------------------
\lstdefinestyle{mmb}{
  basicstyle=\ttfamily\small\color{ink},
  backgroundcolor=\color{panel},
  frame=leftline, framerule=1.2pt, rulecolor=\color{rule},
  framexleftmargin=6pt,
  xleftmargin=10pt, xrightmargin=4pt,
  aboveskip=10pt, belowskip=10pt,
  breaklines=true, columns=fullflexible, keepspaces=true,
  showstringspaces=false,
}
\lstset{style=mmb}

% ---- headings ----------------------------------------------------------
\titleformat{\section}{\Large\bfseries}{\thesection}{0.6em}{}
\titleformat{\subsection}{\large\bfseries}{\thesubsection}{0.6em}{}
\pagestyle{fancy}\fancyhf{}
\fancyhead[L]{\small\textsc{MM-B1 Reference Manual}}
\fancyhead[R]{\small\thepage}
\renewcommand{\headrulewidth}{0.3pt}

\newcommand{\w}[1]{\texttt{#1}}

\begin{document}

% ======================================================================
\begin{titlepage}
\centering
\vspace*{3.5cm}
{\fontsize{64}{64}\selectfont\bfseries MM\,--\,B1}\\[10pt]
{\Large A Decimal Business RPN Calculator}\\[6pt]
{\large exact base-10 money $\cdot$ TVM $\cdot$ cash flows $\cdot$ bonds $\cdot$ dates}\\[28pt]
\rule{0.55\textwidth}{0.4pt}\\[10pt]
{\large User Manual}\\[4pt]
{Version 1.1 \quad$\cdot$\quad \today}\\
\vfill
{\small In the HP-12C lineage; companion to MM-11 and MM-14.\\
All arithmetic via the vendored IBM decNumber reference library. \\
Released under the GNU General Public License, version 3.}
\end{titlepage}

\tableofcontents
\newpage

% ======================================================================
\section{Introduction}

MM-B1 is a reverse-Polish (RPN) calculator for financial work. It belongs to
the same family as the author's earlier MM-11 and MM-14 calculators --- the
same tokenizer, stack, registers, and word-dispatch loop on GNU readline ---
but it differs in two ways that matter for money.

First, \textbf{the arithmetic is decimal}. Binary floating point cannot
represent a value as ordinary as \w{0.10}, so a long chain of binary operations
drifts away from the ledger in the least significant places. MM-B1 computes in
base 10 with explicit round-half-even (banker's) rounding, so its answers agree
with SQL \w{NUMERIC}, COBOL packed decimal, Java \w{BigDecimal}, and Python's
\w{decimal} module --- all implementations of the same General Decimal
Arithmetic Specification. The engine is IBM's decNumber, the C reference
implementation of that specification, vendored with the source so there is no
external decimal dependency.

\begin{lstlisting}
mmb> 0.1 0.2 +
   x |               0.30
\end{lstlisting}

Second, there is a \textbf{real calendar} underneath the bond and cash-flow
functions: actual day counts, the common day-count conventions, accrued
interest, and settlement dates, so a bond can be priced on any business day
rather than only on a coupon date.

\subsection{Conventions used in this manual}

Calculator sessions are shown in panels. What you type follows the \w{mmb>}
prompt; the lines below it are what the calculator prints. The live display
labels the stack levels \w{x}, \w{y}, \w{z}, \w{t} from the bottom up, with
\w{x} the most recent value. In running text, words are shown like \w{?pmt}.

% ======================================================================
\section{Installing and building}

The only external dependency is GNU readline; the decimal library ships with
the source.

\begin{lstlisting}
make            build ./mmb1
make test       run the regression script tests.mmb
make run        build and start an interactive session
\end{lstlisting}

On Linux install readline first (\w{sudo apt install libreadline-dev}). On
macOS readline is keg-only under Homebrew, so point the build at it:

\begin{lstlisting}
brew install readline
make READLINE_PREFIX="$(brew --prefix readline)"
\end{lstlisting}

The decimal library is portable C and compiles identically on arm64 and
x86\_64.

% ======================================================================
\section{The stack and RPN basics}

In RPN you enter operands first and the operation last: \w{3 4 +} leaves
\w{7}. There is no equals key and there are no parentheses; the stack holds
intermediate results. After every line MM-B1 prints the whole stack, newest at
the bottom --- you never need a print command interactively.

\begin{lstlisting}
mmb> 3 4 +
   x |               7.00
mmb> 10 *
   x |              70.00
mmb> 1234567.5 dup +
   y |              70.00
   x |       2,469,135.00
\end{lstlisting}

The bottom four levels are named \w{x y z t}; deeper levels are numbered. Use
\w{dup}, \w{drop}, \w{swap}, and \w{over} to rearrange the stack, \w{clst} to
clear it, and \w{depth} to push the current count. \w{N fix} sets the number of
displayed decimal places; the stored value always keeps full 34-digit
precision. Arithmetic words include \w{+ - * /}, power \w{\^}, \w{chs}
(negate), \w{inv} (reciprocal), \w{sqrt}, \w{ln}, \w{exp}, and \w{abs}.

A line beginning text after \w{\#} is a comment. When input is piped from a
file the stack is \emph{not} auto-printed --- use \w{.} to print and pop the
top value, or \w{.s} to show the stack --- so scripts stay clean and
predictable.

% ======================================================================
\section{Time value of money}

MM-B1 uses a \emph{periodic} TVM model. The rate \w{i} is the interest rate
\textbf{per period, in percent}, and \w{n} is the \textbf{number of periods}.
A loan at 6\% per year paid monthly is entered as \w{0.5 i} with \w{n} counted
in months. There is no separate ``payments per year'' setting to reconcile.

The five registers satisfy
\[
\mathit{PV}\,(1+r)^n
+ \mathit{PMT}\,(1+r\,\tau)\,\frac{(1+r)^n-1}{r}
+ \mathit{FV} = 0, \qquad r=\frac{i}{100},
\]
where $\tau=1$ in \w{begin} mode (annuity due) and $0$ in \w{end} mode (the
default). The \textbf{sign convention} is that cash received is positive and
cash paid is negative. A solved payment that comes back negative is telling you
it is an outflow.

Set any four of \w{n i pv pmt fv} (each consumes the top of the stack) and
solve for the fifth with \w{?n}, \w{?i}, \w{?pv}, \w{?pmt}, or \w{?fv}; the
solver stores its answer back into the register and pushes it. \w{tvm} prints
the register state, \w{clrtvm} resets it.

\subsection{Example: a mortgage}

A \$300{,}000 loan over 30 years at 6\% annual interest, paid monthly.

\begin{lstlisting}
clrtvm
360 n  0.5 i  300000 pv  0 fv  end
?pmt          ->  -1,798.65     monthly payment
\end{lstlisting}

\w{amort} then prints the full schedule --- interest, principal, and remaining
balance for each period --- from the current registers.

\subsection{Example: saving toward a goal}

Depositing \$500 at the end of each month for 30 years at 6\% annual:

\begin{lstlisting}
clrtvm
360 n  0.5 i  0 pv  -500 pmt
?fv           ->  502,257.52    balance after 30 years
\end{lstlisting}

To reach \$1{,}000{,}000 in 25 years at 7\% annual, depositing at the
\emph{start} of each month, switch to annuity due:

\begin{lstlisting}
clrtvm
300 n  0.5833333333 i  0 pv  1000000 fv  begin
?pmt          ->  -1,227.30     required monthly deposit
\end{lstlisting}

\subsection{Example: the true rate on a loan}

Borrow \$10{,}000, repay \$500 a month for 24 months:

\begin{lstlisting}
clrtvm
24 n  10000 pv  -500 pmt  0 fv
?i            ->  1.5131        percent per month (about 18.16% APR)
\end{lstlisting}

% ======================================================================
\section{Cash-flow analysis}

For a stream of cash flows at regular periods, build the cash-flow bank and
evaluate net present value or internal rate of return. \w{cf} appends a flow
(the first one sits at period 0 and is not discounted); \w{cfn} sets a repeat
count on the most recently entered flow; \w{cfs} lists the bank and \w{clrcf}
empties it. \w{npv} takes a periodic rate in percent; \w{irr} takes nothing and
returns the rate in percent.

\begin{lstlisting}
clrcf
-50000 cf  15000 cf  5 cfn
10 npv        ->  6,861.80      net present value at 10%
irr           ->  15.2382       internal rate of return, % per period
\end{lstlisting}

% ======================================================================
\section{Dated cash flows: XNPV and XIRR}

When flows fall on actual calendar dates rather than even periods, use the
dated bank. A date is the plain integer \w{YYYYMMDD}. Append \w{amount date}
pairs with \w{xcf}; \w{xcfs} lists them and \w{clrxcf} clears them. \w{xnpv}
takes an annual rate in percent and discounts each flow by actual/365 from the
date of the first flow --- the same convention as the spreadsheet XNPV and
XIRR functions. \w{xirr} returns the annual rate that sets XNPV to zero.

\begin{lstlisting}
clrxcf
-10000 20080101 xcf
2750 20080301 xcf  4250 20081030 xcf  3250 20090215 xcf  2750 20090401 xcf
9 xnpv        ->  2,086.65      value at 9% per year
xirr          ->  37.3363       money-weighted annual return, %
\end{lstlisting}

% ======================================================================
\section{Interest-rate conversions}

\w{nom>eff} converts a nominal annual rate compounded $m$ times a year into the
effective annual rate; \w{eff>nom} does the reverse. Both take the rate in
percent and $m$ as the number of compounding periods.

\begin{lstlisting}
18 12 nom>eff ->  19.5618       effective annual rate of an 18% APR card
\end{lstlisting}

% ======================================================================
\section{Depreciation}

Three methods are provided, each returning the depreciation for the requested
year. \w{sl} (straight line) is constant and needs only cost, salvage, and
life. \w{soyd} (sum-of-years-digits) and \w{db} (declining balance) also take
the year; \w{db} additionally takes a factor (2 for double-declining balance).

\begin{lstlisting}
50000 5000 7 sl          ->   6,428.57    straight line (every year)
50000 5000 7 1 soyd      ->  11,250.00    sum-of-years-digits, year 1
50000 5000 7 1 2 db      ->  14,285.71    double-declining, year 1
\end{lstlisting}

% ======================================================================
\section{Bonds}

\subsection{Periodic pricing}

\w{bprice} and \w{byield} price a bond on a coupon date with no calendar: a
periodic coupon and yield (in percent), the number of whole periods, and the
redemption value per 100 face.

\begin{lstlisting}
3 3.5 20 100 bprice      ->  92.89        coupon below yield, so a discount
92.89 3 20 100 byield    ->   3.5003      yield recovered from the price
\end{lstlisting}

\subsection{Dated pricing with accrued interest}

\w{xprice} and \w{xyield} price a bond for settlement on any date. They take
the settlement and maturity dates, the annual coupon and yield in percent, the
redemption per 100, the coupon frequency (coupons per year), and the day-count
basis. The basis follows the spreadsheet numbering: 0 is 30/360 (US), 1 is
actual/actual, 2 is actual/360, 3 is actual/365, and 4 is 30E/360 (European).
Coupon dates are generated backward from maturity with end-of-month handling,
and the result is the clean price --- accrued interest netted out. \w{xaccr}
returns the accrued interest by itself.

\begin{lstlisting}
# settle    maturity  cpn% yld% redeem freq basis
20260615 20360515 4.25 4.5 100 2 1 xprice   ->  98.0146   clean price
20260615 20360515 4.25 2 1        xaccr     ->   0.3580   accrued per 100
\end{lstlisting}

The dirty (invoice) price is the clean price plus accrued interest. \w{xyield}
inverts \w{xprice}: give it a price in place of a yield and it returns the
yield. These functions reproduce the documented spreadsheet PRICE and YIELD
worked examples to the displayed precision.

% ======================================================================
\section{Dates and day counts}

A date is the integer \w{YYYYMMDD}; for example \w{20260315} is 15 March 2026.
\w{ddays} returns the actual number of days between two dates, \w{date+} shifts
a date by a number of days, and \w{dow} returns the day of week (0 = Sunday).
\w{yearfrac} returns the fraction of a year between two dates under a chosen
basis (the same 0--4 numbering as the bonds).

\begin{lstlisting}
20070115 20070315 ddays        ->  59           actual days between
20260315 30 date+              ->  20260414     shift a date by days
20260315 dow                   ->  0            day of week (0=Sun..6=Sat)
20070101 20080101 3 yearfrac   ->  1.0000       year fraction, act/365
\end{lstlisting}

% ======================================================================
\section{Numerical methods and conventions}

\begin{itemize}[leftmargin=1.4em]
\item Working precision is 34 significant digits with round-half-even. \w{fix}
      changes only the display (0--16 places); stored values are never
      truncated.
\item The rate solvers (\w{?i}, \w{irr}, \w{xirr}, \w{byield}, \w{xyield}) have
      no closed form. MM-B1 scans the rate axis for a sign change and then
      bisects in decimal to convergence. Where a cash-flow stream admits more
      than one internal rate of return, only the first bracket found is
      reported.
\item \w{amort} assumes end-of-period payments and works in magnitudes. Round
      the payment (or solve \w{?pmt} at the display precision) if you want the
      final balance to close to exactly zero.
\item \w{xnpv} and \w{xirr} fix the year length at 365 days, matching the
      spreadsheet convention; \w{yearfrac} basis 1 uses ISDA actual/actual. For
      dated bonds, basis 0 (30/360) and basis 1 (actual/actual) are the
      best-tested paths.
\end{itemize}

% ======================================================================
\section{Word reference}

\renewcommand{\arraystretch}{1.25}
\noindent
\begin{tabular}{@{}p{0.30\textwidth}p{0.64\textwidth}@{}}
\toprule
\textbf{Group} & \textbf{Words} \\
\midrule
Arithmetic        & \w{+ - * / \^{} chs inv sqrt ln exp abs} \\
Business percent  & \w{pct} ($x$\% of $y$), \w{\%chg}, \w{markup}, \w{margin} \\
Stack             & \w{dup drop swap over depth clst . .s} \\
Constants         & \w{pi e} \\
TVM set           & \w{n i pv pmt fv} \\
TVM solve         & \w{?n ?i ?pv ?pmt ?fv} \\
TVM mode          & \w{begin end clrtvm tvm} \\
Cash flows        & \w{cf cfn clrcf cfs npv irr} \\
Amortization      & \w{amort} \\
Rate conversion   & \w{nom>eff eff>nom} \\
Depreciation      & \w{sl soyd db} \\
Bonds (periodic)  & \w{bprice byield} \\
Dates             & \w{ddays date+ dow yearfrac} \\
Dated cash flows  & \w{xcf clrxcf xcfs xnpv xirr} \\
Dated bonds       & \w{xprice xyield xaccr} \\
Storage           & \w{s0..s9} (store, leaves stack), \w{r0..r9} (recall) \\
Display / meta    & \w{fix help words quit} \\
\bottomrule
\end{tabular}

\vspace{1em}
\noindent The argument order is always postfix. A few of the multi-argument
words, in the order their operands are pushed:

\begin{lstlisting}
y x pct                                          x% of y
cost salv life sl                                straight-line depreciation
cost salv life year factor db                    declining-balance depreciation
nom% m nom>eff                                    nominal -> effective annual
cpn% yld% n redeem bprice                         periodic bond price
settle mat cpn% yld% redeem freq basis xprice     dated bond clean price
settle mat cpn% price redeem freq basis xyield    dated bond yield
settle mat cpn% freq basis xaccr                  accrued interest per 100
amount date xcf                                   append a dated cash flow
\end{lstlisting}

% ======================================================================
\section{License}

MM-B1 is released under the GNU General Public License, version 3, consistent
with MM-11 and MM-14. The vendored decNumber library is distributed under the
permissive ICU license, which the GPL permits redistributing alongside.

\end{document}
